% Appendix C (rewritten 2026-10-16 by the appendix writer, <= 1 page). Sources: the pre-condense design section's "The
% external harness and the equal footprint" and "Thermal governance as part of the measurement" (fact-checked
% 2026-10-16), D-003, D-008, D-009, D-021; the footprint table as in the C1 re-read's run records. table_calib and
% table_chromatic are generated by tools/paper_tables.py.
\section{Measurement harness and calibration}\label{app:harness}
Every compared GPU arm gets the same machine share $F$ (Table~\ref{tab:footprint}). Thermal rules are part of the
measurement, since a capped GPU would bias a comparison; the runner enforces them (D-003, D-008, D-009, D-021). The
120-s budget starts at a shared barrier, after every process has read the instance and warmed its GPU kernels.

\begin{table}[!tb]\centering\scriptsize
\caption{The machine share $F$ under the final harness. Memory caps are hard (systemd \texttt{MemoryMax} on the
host, a pool limit on the GPU); the runner's own threads run outside $F$.}\label{tab:footprint}
\setlength{\tabcolsep}{3pt}
\begin{tabular}{>{\raggedright}p{4.3cm}>{\raggedright}p{3.6cm}>{\raggedright\arraybackslash}p{7.6cm}}
\toprule
process & cores & memory caps and settings \\
\midrule
SCIP~10, one thread & Cortex-X925 core 5 & host 9\,GiB (SCIP's own limit 8\,GiB) \\
EA and its two HiGHS processes & Cortex-X925 core 6; GPU & host 6\,GiB; GPU pool 12\,GiB; CUDA yield flag \\
cuOpt~26.08, three threads & Cortex-A725 cores 1--3; GPU & host and GPU 24\,GiB each; MIP gap limits 0; thread
  variables (OMP, OpenBLAS, MKL) unset; own kernel cache \\
cuOpt alone, five threads & all five cores; GPU & as cuOpt above \\
\bottomrule
\end{tabular}
\end{table}

\paragraph{Thermal rules.} The host's fan governor caps the GPU clock at 2{,}200\,MHz after 20\,s at or above
85\,$^\circ$C at maximum fan, and requests a 30-minute rest after 30 minutes at maximum fan.
\begin{itemize}\setlength{\itemsep}{0pt}\setlength{\parskip}{0pt}
\item \textbf{Windows and start gate:} daytime windows under a lease, with the local language-model servers stopped.
  A GPU run starts below 80\,$^\circ$C (CPU-only: 88\,$^\circ$C), with no rest requested and no clock cap; rests took
  about a quarter of window time.
\item \textbf{Voids:} a block (one instance and seed, both arms) is voided, requeued once and rerun whole if the hottest zone
  reaches 92\,$^\circ$C, a GPU run's median busy SM clock is below 2{,}300\,MHz or meets a clock cap, another program's GPU
  process runs for over 2\,\% of the budget, or a kernel compiles on the clock. Below 2{,}000\,MHz (a known stuck-clock
  fault) the run set halts.
\item \textbf{Breaker:} two thermal or cap voids close the window. In the C1 re-read it fired once (15 October, D-030);
  the remaining blocks ran whole the next morning.
\end{itemize}

\paragraph{Harness measurements} (D-021), made before any run on unseen instances; they test the harness, not the
method. M4 reads outside its band, which D-021 attributes to $F$ (five threads and the GPU, against a 24-thread whole
machine) and to the machine, so the unseen-instance test's targets are anchored on our own cuOpt-alone runs.
\gentable{list_calib}
Table~\ref{tab:chromatic} gives the per-run values behind Table~\ref{tab:loss}.

\gentable{table_chromatic}
